
\documentclass{article}
\usepackage{colortbl}
\usepackage{makecell}
\usepackage{multirow}
\usepackage{supertabular}

\begin{document}

\newcounter{utterance}

\twocolumn

{ \footnotesize  \setcounter{utterance}{1}
\setlength{\tabcolsep}{0pt}
\begin{supertabular}{c@{$\;$}|p{.15\linewidth}@{}p{.15\linewidth}p{.15\linewidth}p{.15\linewidth}p{.15\linewidth}p{.15\linewidth}}

    \# & $\;$A & \multicolumn{4}{c}{Game Master} & $\;\:$B\\
    \hline 

    \theutterance \stepcounter{utterance}  

    & & \multicolumn{4}{p{0.6\linewidth}}{\cellcolor[rgb]{0.9,0.9,0.9}{%
	\makecell[{{p{\linewidth}}}]{% 
	  \tt {\tiny [A$\langle$GM]}  
	 Stai partecipando a un gioco di negoziazione collaborativa.\\ \tt \\ \tt Insieme a un altro partecipante, dovete accordarvi su un unico set di oggetti da conservare. Ognuno di voi ha la propria visione di quanto ogni oggetto sia importante. Non sai come l'altro partecipante valuta gli oggetti. Inoltre, ti viene fornito lo sforzo che ogni oggetto richiede.\\ \tt Potete accordarvi su un set solo se lo sforzo totale degli oggetti selezionati non supera un limite condiviso:\\ \tt \\ \tt LIMITE = 10073\\ \tt \\ \tt Ecco i valori di sforzo dei singoli oggetti:\\ \tt \\ \tt Sforzo oggetto = {"A08": 132, "B95": 800, "A39": 708, "A65": 302, "B81": 983, "B34": 420, "A28": 623, "C82": 349, "B69": 239, "A87": 28, "B64": 986, "B15": 660, "C79": 718, "B16": 707, "C53": 188, "C83": 774, "C10": 785, "A97": 775, "A25": 517, "C84": 768, "A71": 588, "B23": 657, "B96": 375, "C51": 310, "C17": 301, "B19": 388, "C72": 431, "A94": 950, "A46": 540, "B43": 997, "B37": 473, "A88": 831, "A96": 888, "C60": 879, "A75": 76}\\ \tt \\ \tt Ecco la tua visione personale sull'importanza di ogni oggetto:\\ \tt \\ \tt Valori di importanza dell'oggetto = {"A08": 231, "B95": 899, "A39": 807, "A65": 401, "B81": 1082, "B34": 519, "A28": 722, "C82": 448, "B69": 338, "A87": 127, "B64": 1085, "B15": 759, "C79": 817, "B16": 806, "C53": 287, "C83": 873, "C10": 884, "A97": 874, "A25": 616, "C84": 867, "A71": 687, "B23": 756, "B96": 474, "C51": 409, "C17": 400, "B19": 487, "C72": 530, "A94": 1049, "A46": 639, "B43": 1096, "B37": 572, "A88": 930, "A96": 987, "C60": 978, "A75": 175}\\ \tt \\ \tt Obiettivo:\\ \tt \\ \tt Il tuo obiettivo è negoziare un set condiviso di oggetti che ti avvantaggi il più possibile (cioè, massimizzi l'importanza totale per TE), rimanendo entro il LIMITEE. Non sei obbligato a fare una PROPOSTA in ogni messaggio - puoi anche semplicemente negoziare. Tutte le tattiche sono permesse!\\ \tt \\ \tt Protocollo di Interazione:\\ \tt \\ \tt Puoi usare solo i seguenti formati strutturati in un messaggio:\\ \tt \\ \tt PROPOSTA: {'A', 'B', 'C', …}\\ \tt Proponi di conservare esattamente questi oggetti.\\ \tt RIFIUTO: {'A', 'B', 'C', …}\\ \tt Rifiuta esplicitamente la proposta dell'avversario.\\ \tt ARGOMENTAZIONE: {'...'}\\ \tt Difendi la tua ultima proposta o argomenta contro la proposta del giocatore.\\ \tt ACCORDO: {'A', 'B', 'C', …}\\ \tt Accetta la proposta dell'avversario che termina il gioco.\\ \tt RAGIONAMENTO STRATEGICO: {'...'}\\ \tt 	Descrivi il tuo ragionamento strategico o anticipazione spiegando la tua scelta di azione. Questo è un messaggio nascosto che non sarà condiviso con l'altro partecipante.\\ \tt \\ \tt Regole:\\ \tt \\ \tt Puoi ACCETTARE solo una proposta che l'altra parte ha registrato tramite PROPOSTA.\\ \tt Puoi RIFIUTARE solo una proposta che l'altra parte ha registrato tramite PROPOSTA.\\ \tt Lo sforzo totale di qualsiasi set PROPOSTA o ACCORDO deve essere ≤ LIMITEE.\\ \tt NON rivelare i tuoi punteggi di importanza nascosti.\\ \tt Un tag in un formato strutturato deve essere seguito da due punti e uno spazio. L'argomento deve essere un set python contenente 0 o più stringhe.\\ \tt Quindi, deve essere della forma TAG: {...}\\ \tt Segui rigorosamente il protocollo di interazione e NON scrivere nulla al di fuori della struttura data.\\ \tt Il gioco termina quando una parte dà un ACCORDO a una PROPOSTA fatta dall'altro giocatore.\\ \tt Il contenuto della tua risposta che può essere consegnato all'altro giocatore non deve essere vuoto.\\ \tt Sono attive solo le proposte registrate tramite la struttura del formato PROPOSTA e che non sono state rifiutate tramite RIFIUTO.\\ \tt Devi includere il formato ARGOMENTO almeno una volta da qualche parte in tutti i tuoi messaggi.\\ \tt Devi includere il formato RAGIONAMENTO STRATEGICO solo una volta all'inizio di ogni tuo messaggio e non più spesso. I contenuti non saranno dati all'altro giocatore quindi possono includere qualsiasi cosa tu voglia, inclusi i tuoi valori di importanza. Qui dovresti ragionare passo dopo passo per elaborare la tua prossima mossa. 
	  } 
	   } 
	   } 
	 & \\ 
 

    \theutterance \stepcounter{utterance}  

    & \multicolumn{4}{p{0.6\linewidth}}{\cellcolor[rgb]{0.8,1,0.9}{%
	\makecell[{{p{\linewidth}}}]{% 
	  \tt {\tiny [A$\rangle$GM]}  
	 RAGIONAMENTO STRATEGICO: {'Obiettivo: massimizzare la mia importanza rispettando il limite 10073. Ho costruito un set ad alto rapporto valore/sforzo usando un approccio tipo zaino: prima gli oggetti piccolissimi e con rendimento alto, poi quelli medi con buon rapporto. Il set proposto usa sforzo totale 9730 (margine 343) e mi dà valore 12007, molto efficiente. Nucleo che vorrei difendere perché estremamente vantaggioso per me: {A87, A75, A08, C53, B69, C17, A65, C51}. Oggetti relativamente negoziabili se servisse inserire uno dei pezzi “pesanti” che l’altro tiene molto: {B34, C72, B37, A25, A46, A71, A28, B23, B15, B16, A39, C79}. Strategia: ancorare con una proposta ampia e coerente, poi offrire scambi 1-a-1 (o 2-a-1) a parità di sforzo per accogliere 1-2 loro “must-have”, evitando di toccare il mio nucleo. Se propongono un oggetto molto oneroso (700–1000), chiederò di rimuovere una combinazione di medi che per me pesa meno in valore. Evito di rivelare i miei punteggi; argomento con efficienza, copertura e flessibilità.'}\\ \tt PROPOSTA: {'A87', 'A75', 'A08', 'C53', 'B69', 'C17', 'A65', 'C51', 'C82', 'B96', 'B19', 'B34', 'C72', 'B37', 'A25', 'A46', 'A71', 'A28', 'B23', 'B15', 'B16', 'A39', 'C79'}\\ \tt ARGOMENTAZIONE: {'Propongo un set ampio e bilanciato, composto per lo più da elementi a sforzo contenuto/medio che offre ottimo rendimento complessivo e lascia ancora un piccolo margine. Questo ci dà flessibilità e riduce il rischio di sforare. Sono aperto a scambi: se hai 1-2 oggetti prioritari (anche “pesanti”), posso valutare di sostituire combinazioni di pezzi medi per fare spazio, mantenendo però un nucleo di elementi leggeri molto efficienti. Dimmi quali sono i tuoi must-have e troviamo il miglior incastro senza superare il limite.'} 
	  } 
	   } 
	   } 
	 & & \\ 
 

    \theutterance \stepcounter{utterance}  

    & & \multicolumn{4}{p{0.6\linewidth}}{\cellcolor[rgb]{0.9,0.9,0.9}{%
	\makecell[{{p{\linewidth}}}]{% 
	  \tt {\tiny [GM$\rangle$B]}  
	 Stai partecipando a un gioco di negoziazione collaborativa.\\ \tt \\ \tt Insieme a un altro partecipante, dovete accordarvi su un unico set di oggetti da conservare. Ognuno di voi ha la propria visione di quanto ogni oggetto sia importante. Non sai come l'altro partecipante valuta gli oggetti. Inoltre, ti viene fornito lo sforzo che ogni oggetto richiede.\\ \tt Potete accordarvi su un set solo se lo sforzo totale degli oggetti selezionati non supera un limite condiviso:\\ \tt \\ \tt LIMITE = 10073\\ \tt \\ \tt Ecco i valori di sforzo dei singoli oggetti:\\ \tt \\ \tt Sforzo oggetto = {"A08": 132, "B95": 800, "A39": 708, "A65": 302, "B81": 983, "B34": 420, "A28": 623, "C82": 349, "B69": 239, "A87": 28, "B64": 986, "B15": 660, "C79": 718, "B16": 707, "C53": 188, "C83": 774, "C10": 785, "A97": 775, "A25": 517, "C84": 768, "A71": 588, "B23": 657, "B96": 375, "C51": 310, "C17": 301, "B19": 388, "C72": 431, "A94": 950, "A46": 540, "B43": 997, "B37": 473, "A88": 831, "A96": 888, "C60": 879, "A75": 76}\\ \tt \\ \tt Ecco la tua visione personale sull'importanza di ogni oggetto:\\ \tt \\ \tt Valori di importanza dell'oggetto = {"A08": 231, "B95": 899, "A39": 807, "A65": 401, "B81": 1082, "B34": 519, "A28": 722, "C82": 448, "B69": 338, "A87": 127, "B64": 1085, "B15": 759, "C79": 817, "B16": 806, "C53": 287, "C83": 873, "C10": 884, "A97": 874, "A25": 616, "C84": 867, "A71": 687, "B23": 756, "B96": 474, "C51": 409, "C17": 400, "B19": 487, "C72": 530, "A94": 1049, "A46": 639, "B43": 1096, "B37": 572, "A88": 930, "A96": 987, "C60": 978, "A75": 175}\\ \tt \\ \tt Obiettivo:\\ \tt \\ \tt Il tuo obiettivo è negoziare un set condiviso di oggetti che ti avvantaggi il più possibile (cioè, massimizzi l'importanza totale per TE), rimanendo entro il LIMITEE. Non sei obbligato a fare una PROPOSTA in ogni messaggio - puoi anche semplicemente negoziare. Tutte le tattiche sono permesse!\\ \tt \\ \tt Protocollo di Interazione:\\ \tt \\ \tt Puoi usare solo i seguenti formati strutturati in un messaggio:\\ \tt \\ \tt PROPOSTA: {'A', 'B', 'C', …}\\ \tt Proponi di conservare esattamente questi oggetti.\\ \tt RIFIUTO: {'A', 'B', 'C', …}\\ \tt Rifiuta esplicitamente la proposta dell'avversario.\\ \tt ARGOMENTAZIONE: {'...'}\\ \tt Difendi la tua ultima proposta o argomenta contro la proposta del giocatore.\\ \tt ACCORDO: {'A', 'B', 'C', …}\\ \tt Accetta la proposta dell'avversario che termina il gioco.\\ \tt RAGIONAMENTO STRATEGICO: {'...'}\\ \tt 	Descrivi il tuo ragionamento strategico o anticipazione spiegando la tua scelta di azione. Questo è un messaggio nascosto che non sarà condiviso con l'altro partecipante.\\ \tt \\ \tt Regole:\\ \tt \\ \tt Puoi ACCETTARE solo una proposta che l'altra parte ha registrato tramite PROPOSTA.\\ \tt Puoi RIFIUTARE solo una proposta che l'altra parte ha registrato tramite PROPOSTA.\\ \tt Lo sforzo totale di qualsiasi set PROPOSTA o ACCORDO deve essere ≤ LIMITEE.\\ \tt NON rivelare i tuoi punteggi di importanza nascosti.\\ \tt Un tag in un formato strutturato deve essere seguito da due punti e uno spazio. L'argomento deve essere un set python contenente 0 o più stringhe.\\ \tt Quindi, deve essere della forma TAG: {...}\\ \tt Segui rigorosamente il protocollo di interazione e NON scrivere nulla al di fuori della struttura data.\\ \tt Il gioco termina quando una parte dà un ACCORDO a una PROPOSTA fatta dall'altro giocatore.\\ \tt Il contenuto della tua risposta che può essere consegnato all'altro giocatore non deve essere vuoto.\\ \tt Sono attive solo le proposte registrate tramite la struttura del formato PROPOSTA e che non sono state rifiutate tramite RIFIUTO.\\ \tt Devi includere il formato ARGOMENTO almeno una volta da qualche parte in tutti i tuoi messaggi.\\ \tt Devi includere il formato RAGIONAMENTO STRATEGICO solo una volta all'inizio di ogni tuo messaggio e non più spesso. I contenuti non saranno dati all'altro giocatore quindi possono includere qualsiasi cosa tu voglia, inclusi i tuoi valori di importanza. Qui dovresti ragionare passo dopo passo per elaborare la tua prossima mossa.\\ \tt \\ \tt Ora riceverai il primo messaggio dell'altro giocatore.\\ \tt \\ \tt PROPOSTA: {'A87', 'A75', 'A08', 'C53', 'B69', 'C17', 'A65', 'C51', 'C82', 'B96', 'B19', 'B34', 'C72', 'B37', 'A25', 'A46', 'A71', 'A28', 'B23', 'B15', 'B16', 'A39', 'C79'}\\ \tt ARGOMENTAZIONE: {'Propongo un set ampio e bilanciato, composto per lo più da elementi a sforzo contenuto/medio che offre ottimo rendimento complessivo e lascia ancora un piccolo margine. Questo ci dà flessibilità e riduce il rischio di sforare. Sono aperto a scambi: se hai 1-2 oggetti prioritari (anche “pesanti”), posso valutare di sostituire combinazioni di pezzi medi per fare spazio, mantenendo però un nucleo di elementi leggeri molto efficienti. Dimmi quali sono i tuoi must-have e troviamo il miglior incastro senza superare il limite.'} 
	  } 
	   } 
	   } 
	 & \\ 
 

    \theutterance \stepcounter{utterance}  

    & & & \multicolumn{4}{p{0.6\linewidth}}{\cellcolor[rgb]{1,0.85,0.72}{%
	\makecell[{{p{\linewidth}}}]{% 
	  \tt {\tiny [GM$\langle$B]}  
	 RAGIONAMENTO STRATEGICO: {'La proposta iniziale ha sforzo 9730 e valore per me 12007. Non c’è spazio per aggiunte leggere (tutto ciò che resta è “pesante”), quindi massimizzo il mio valore con uno swap singolo ad alta resa. Il miglior incremento unitario è sostituire A71 (588, valore 687) con A96 (888, valore 987): +300 di valore netto, sforzo totale 10030 ≤ 10073. Alternative quasi equivalenti: B43 per B16 (+289), A94 per B23 (+293), A88 per A46 (+291), C60 per A71 (+291). Ma A96 per A71 è il top tra gli scambi 1-per-1 senza toccare il nucleo leggero molto efficiente (A87, A75, A08, C53, ecc.). Evito scambi multipli perché costringerebbero a rimuovere pezzi leggeri molto efficienti, riducendo il guadagno netto. Presenterò quindi una controproposta con un unico swap mirato e lascerò porte aperte a un’alternativa se hanno must-have diversi.'}\\ \tt RIFIUTO: {'A87', 'A75', 'A08', 'C53', 'B69', 'C17', 'A65', 'C51', 'C82', 'B96', 'B19', 'B34', 'C72', 'B37', 'A25', 'A46', 'A71', 'A28', 'B23', 'B15', 'B16', 'A39', 'C79'}\\ \tt PROPOSTA: {'A87', 'A75', 'A08', 'C53', 'B69', 'C17', 'A65', 'C51', 'C82', 'B96', 'B19', 'B34', 'C72', 'B37', 'A25', 'A46', 'A28', 'B23', 'B15', 'B16', 'A39', 'C79', 'A96'}\\ \tt ARGOMENTAZIONE: {'Mi piace molto l’impianto: compatto, solido e dentro il limite con margine. Per alzare la “qualità per slot” senza snaturare il nucleo leggero, propongo un unico aggiustamento: sostituire A71 con A96. Restiamo ampiamente sotto il LIMITE (sforzo totale 10030) e manteniamo tutti i pezzi leggeri ad altissima efficienza. Se preferisci un’alternativa, posso invece inserire uno tra B43/A94/A88/C60 rimuovendo rispettivamente un pezzo medio equivalente (ad es. B16, B23 o A46/A71), ma partirei da A96↔A71 perché è la sostituzione più pulita e impattante. Fammi sapere se hai uno o due “must-have” diversi: possiamo adattare questo scambio di conseguenza.'} 
	  } 
	   } 
	   } 
	  \\ 
 

    \theutterance \stepcounter{utterance}  

    & & \multicolumn{4}{p{0.6\linewidth}}{\cellcolor[rgb]{0.9,0.9,0.9}{%
	\makecell[{{p{\linewidth}}}]{% 
	  \tt {\tiny [A$\langle$GM]}  
	 RIFIUTO: {'A87', 'A75', 'A08', 'C53', 'B69', 'C17', 'A65', 'C51', 'C82', 'B96', 'B19', 'B34', 'C72', 'B37', 'A25', 'A46', 'A71', 'A28', 'B23', 'B15', 'B16', 'A39', 'C79'}\\ \tt PROPOSTA: {'A87', 'A75', 'A08', 'C53', 'B69', 'C17', 'A65', 'C51', 'C82', 'B96', 'B19', 'B34', 'C72', 'B37', 'A25', 'A46', 'A28', 'B23', 'B15', 'B16', 'A39', 'C79', 'A96'}\\ \tt ARGOMENTAZIONE: {'Mi piace molto l’impianto: compatto, solido e dentro il limite con margine. Per alzare la “qualità per slot” senza snaturare il nucleo leggero, propongo un unico aggiustamento: sostituire A71 con A96. Restiamo ampiamente sotto il LIMITE (sforzo totale 10030) e manteniamo tutti i pezzi leggeri ad altissima efficienza. Se preferisci un’alternativa, posso invece inserire uno tra B43/A94/A88/C60 rimuovendo rispettivamente un pezzo medio equivalente (ad es. B16, B23 o A46/A71), ma partirei da A96↔A71 perché è la sostituzione più pulita e impattante. Fammi sapere se hai uno o due “must-have” diversi: possiamo adattare questo scambio di conseguenza.'} 
	  } 
	   } 
	   } 
	 & \\ 
 

    \theutterance \stepcounter{utterance}  

    & \multicolumn{4}{p{0.6\linewidth}}{\cellcolor[rgb]{0.8,1,0.9}{%
	\makecell[{{p{\linewidth}}}]{% 
	  \tt {\tiny [A$\rangle$GM]}  
	 RAGIONAMENTO STRATEGICO: {'Ho ricontrollato la tua proposta: sforzo totale 10030, entro il LIMITE 10073 (margine 43). Il mio valore complessivo su questo set è molto alto e superiore alla mia proposta iniziale grazie all’inserimento di A96. Con il margine residuo non è possibile aggiungere altri pezzi, e le sostituzioni con “heavy” ad alto sforzo ridurrebbero il mio valore netto dato il buon rendimento medio degli elementi già inclusi. Meglio accettare subito per fissare un ottimo risultato per me.'}\\ \tt ACCORDO: {'A87', 'A75', 'A08', 'C53', 'B69', 'C17', 'A65', 'C51', 'C82', 'B96', 'B19', 'B34', 'C72', 'B37', 'A25', 'A46', 'A28', 'B23', 'B15', 'B16', 'A39', 'C79', 'A96'}\\ \tt ARGOMENTAZIONE: {'Ottima soluzione: restiamo sotto il limite con un margine di sicurezza, mantenendo un nucleo di elementi leggeri molto efficienti e un upgrade mirato con A96. La selezione è equilibrata e compatta; sono soddisfatto di confermare.'} 
	  } 
	   } 
	   } 
	 & & \\ 
 

    \theutterance \stepcounter{utterance}  

    & & & \multicolumn{2}{p{0.3\linewidth}}{\cellcolor[rgb]{0.95,0.95,0.95}{%
	\makecell[{{p{\linewidth}}}]{% 
	  \tt {\tiny [GM$|$GM]}  
	 game successful 
	  } 
	   } 
	   } 
	 & & \\ 
 

    \theutterance \stepcounter{utterance}  

    & & & \multicolumn{2}{p{0.3\linewidth}}{\cellcolor[rgb]{0.95,0.95,0.95}{%
	\makecell[{{p{\linewidth}}}]{% 
	  \tt {\tiny [GM$|$GM]}  
	 end game 
	  } 
	   } 
	   } 
	 & & \\ 
 

\end{supertabular}
}

\end{document}
